%!TEX root = ../dokumentation.tex

\chapter{Umsetzung}

\section{Planung und Aufbau}

Für den praktischen Aufbau des Systems werden folgende Komponenten benötigt:

\begin{itemize}
	\item DHT11 - Temperatur- und Luftfeuchtigkeitssensor
	\item ESP32 - Mikrocontroller mit WLAN-Funktion
	\item Breadboard - zur schnellen, lötfreien Prototypenentwicklung
	\item Micro-USB-Kabel - zur Stromversorgung und Programmierung
	\item 3 Jumper-Kabel - zur elektrischen Verbindung zwischen Sensor und Mikrocontroller
\end{itemize}

Der Aufbau erfolgt wie in Abbildung 3.1 dargestellt. 
Der DHT11 wird mit drei Leitungen an den ESP32 angeschlossen:
ein Datenpin, ein 3.3-V-Anschluss sowie eine Masseverbindung. 
Das Micro-USB-Kabel verbindet den ESP32 mit einem Computer und
dient sowohl als Stromquelle als auch zur Übertragung des
Programmcodes. Durch die geringe Leistungsaufnahme des 
Systems reicht die USB-Stromversorgung vollständig aus.
\begin{figure} [H]
	 \centering 
	 \resizebox{.5\textwidth}{!}
	 { \includegraphics{circuit_image.png} } 
	 \vspace{10pt} 
     \caption {Schaltplan DHT11 (eigene Abbildung)}\label{Schaltplan} 
\end{figure}

Zur Programmierung wird Visual Studio Code mit der Erweiterung 
PlatformIO genutzt, die Mikrocontrollerprojekte, 
Bibliothekenverwaltung, Kompilieren und Hochladen auf den ESP32 
vereinfacht. Für eine einfache Weboberfläche zur Anzeige der 
Messdaten kommt zusätzlich Live Server zum Einsatz, um die 
Website lokal zu hosten und in Echtzeit zu aktualisieren.

Der Plan sieht vor, den DHT11 in regelmäßigen Abständen 
auszulesen, die Temperatur- und Feuchtigkeitswerte digital 
zu verarbeiten und über den seriellen Monitor zu prüfen. 
Optional können die Daten über die Weboberfläche dargestellt 
werden, um die Genauigkeit des Sensors für Anwendungen wie 
Defogging-Systeme zu beurteilen.



\section{Durchführung}

\subsection{Messreihe zur Genauigkeit des DHT11 Sensors}
Für die Durchführung der Messungen wurde der Versuchsaufbau 
in einem geschlossenen Raum platziert, um konstante 
Umgebungsbedingungen zu gewährleisten. Dann wurden je fünf Messungen mit zwei verschiedenen DHT11 Sensoren gestartet, 
um sicher zu 
gehen, dass der verwendete Sensor nicht durch einen FEertigungsfehler oder Beschädigung ungenauer ist als andere DHT11 Sensoren.
 Der ESP32 liest 
den DHT11 in einem festen Intervall von zwei Sekunden aus. 
Die Messwerte werden unmittelbar über die serielle Schnittstelle
ausgegeben und auf der vom ESP32 gehosteten Website 
ausgegeben.

\begin{figure}[H]
	\centering
    \begin{subfigure}{.49\textwidth}
        \centering
        \resizebox{\textwidth}{!}{
            \includegraphics{Temperaturmessreihe.png} 
            }
	        \vspace{10pt} 
            \caption {Messreihe Temperatur (eigene Abbildung)}\label{Messreihe Temperatur}
    \end{subfigure}
    \begin{subfigure}{.49\textwidth}
            \centering
        \resizebox{\textwidth}{!}{
             \includegraphics{Luftfeuchtigkeitsmessreihe.png} 
        }
	        \vspace{10pt} 
            \caption {Messreihe Luftfeuchtigkeit (eigene Abbildung)}\label{Messreihe Luftfeuchtigkeit} 
    \end{subfigure}
    \caption{Messungen zur Genauigkeit des DHT11 Sensors}\label{Messungen DHT11 (Eigene Abbildung)}
\end{figure}

Die Standardabweichung der Temperaturmessungen über beide 
Sensoren liegt bei 0,32°C, während die Luftfeuchtigkeit eine 
Standardabweichung von 1,75\% aufweist. Diese Werte deuten darauf 
hin, dass die Temperatursensoren eine hohe Präzision aufweisen, 
während die Feuchtigkeitssensoren geringfügig stärkeren 
Schwankungen unterliegen.

Ein Vergleich der beiden Sensoren zeigt, dass Sensor 2 
systematisch leicht höhere Werte sowohl für Temperatur 
(+0,56°C) als auch für Luftfeuchtigkeit (+0,92\%) liefert. Was jedoch
ebenfalls darauf zurückzuführen sein konnte, dass sowohl die Luftfeuchtigkeit
als auch die Temperatur im Raum gestiegen ist, da sich mehrere Menschen dort befanden.


Insgesamt bestätigen die Daten jedoch eine zufriedenstellende 
Übereinstimmung der Sensoren. Die geringe Streuung der Messwerte 
spricht für eine zuverlässige Messgenauigkeit, die für den Einsatz im präventiven Defogging geeignet ist.
 


\section{Theoretische Implementierung zum Defogging}
Im Folgenden wird beschrieben, wie das entwickelte Sensorsystem theoretisch zur Prävention von 
Scheibenbeschlag eingesetzt werden kann. Grundlage hierfür ist das physikalische Prinzip der 
Kondensation: Beschlag entsteht, wenn die Temperatur der Scheibenoberfläche unter den Taupunkt 
der umgebenden Luft fällt \citeM{i}{Ene_2021}{2}{2}. Zur Bestimmung dieses Zusammenhangs werden 
Taupunkttabellen herangezogen.
\begin{figure}[H]
	 \centering 
	 \resizebox{\textwidth}{!}
	 { \includegraphics{ausschnitt_Taupunkttabelle.png} } 
	 \vspace{10pt} 
	 \caption{Ausschnitt aus der Taupunkttabelle \protect\citeM{d}{tauKor}{1}{1}}\label{Ausschnitt_Taupunkttabelle} 
\end{figure}

Die Taupunkttabelle gibt an, bei welchen Kombinationen aus Lufttemperatur und relativer Luftfeuchtigkeit Kondensation auftritt. 
Um diese Bedingungen zuverlässig zu erkennen, ist neben der Messung der Lufttemperatur und -feuchtigkeit durch den DHT11 ein zusätzlicher 
Sensor erforderlich, der die Temperatur der Windschutzscheibe erfasst.

Ein Algorithmus wertet anschließend alle drei Messgrößen Lufttemperatur, Luftfeuchtigkeit und Scheibentemperatur in regelmäßigen Abständen aus. 
Liegt beispielsweise eine Lufttemperatur von 10 °C bei einer relativen Luftfeuchtigkeit von 70 \% vor, so tritt Kondensation auf, 
sobald die Scheibentemperatur auf 4,79 °C oder darunter absinkt (siehe Abbildung \ref{Ausschnitt_Taupunkttabelle}). Der Algorithmus muss folglich nicht nur den aktuellen Zustand erfassen, sondern insbesondere die Entwicklung der Scheibentemperatur überwachen, um rechtzeitig vor dem Erreichen des Taupunkts reagieren zu können.

Als Gegenmaßnahme wird eine leichte Erwärmung der Scheibe vorgesehen. Diese kann entweder über das \ac{HVAC}-System 
oder, sofern vorhanden, über eine integrierte Scheibenheizung erfolgen. 
Ziel ist es, den Temperaturunterschied zwischen Scheibe und Taupunkt so gering zu 
halten, dass eine Kondensation vermieden wird.

\section{Kritische Betrachtung und Lösungsansätze}

Bei der Betrachtung des vorgestellten Konzepts ergeben sich mehrere Einschränkungen und Herausforderungen. 
Für Fahrzeuge mit Verbrennungsmotor ist das System nur bedingt geeignet, da ein automatisches Heizsystem 
die Fahrzeugbatterie unnötig belasten würde. Da die elektrische Energie hier nicht primär für den Fahrzeugantrieb 
vorgesehen ist, steht der zusätzliche Energieverbrauch in keinem sinnvollen Verhältnis zum Nutzen.

Für Elektrofahrzeuge erscheint der Ansatz grundsätzlich geeigneter, da die benötigte Energie aus dem Hauptakku 
bezogen werden kann. Dennoch zeigt sich auch hier ein zentrales Problem: Der präventive Charakter des Systems führt 
in vielen Situationen zu einem ineffizienten Energieeinsatz. Ist das System dauerhaft aktiv, wird kontinuierlich 
Energie für das Auslesen und Auswerten der Sensordaten sowie für eine potenzielle Scheibenbeheizung verbraucht, 
selbst wenn kein Beschlag droht. Insbesondere bei längeren Standzeiten, etwa über Nacht oder bei selten genutzten 
Fahrzeugen, würde das System unnötig Energie verbrauchen und somit die Reichweite eines Elektrofahrzeugs spürbar 
reduzieren.

\subsection{Lösungsansatz 1: Aktivierung nur beim Ladevorgang}
Ein möglicher Ansatz besteht darin, das System ausschließlich während des Ladevorgangs eines Elektrofahrzeugs zu 
aktivieren. In diesem Fall hätte der zusätzliche Energieverbrauch keinen direkten Einfluss auf die Reichweite. 
Allerdings wird dabei weiterhin Energie eingesetzt, auch wenn eine Defogging-Maßnahme nicht erforderlich ist. Der 
grundsätzliche Nachteil eines ineffizienten Dauerbetriebs bleibt somit bestehen, sodass dieser Ansatz zwar praktikabel,
jedoch nicht nachhaltig ist.


\subsection{Lösungsansatz 2: Intelligentes Defogging}
MEine weitere Möglichkeit ist die Implementierung eines intelligenten Systems, das das Fahrverhalten des Nutzers 
analysiert. Erkennt das System beispielsweise wiederkehrende Fahrzeiten, etwa den täglichen Arbeitsweg, könnte das 
Defogging automatisch zu diesen Zeiten aktiviert werden, ohne dass ein manuelles Eingreifen erforderlich ist. Dieser 
Ansatz reduziert zwar unnötigen Energieverbrauch, birgt jedoch ebenfalls Risiken. Ändert sich das Fahrverhalten, etwa 
durch Urlaub oder flexible Arbeitszeiten, kann das System unnötig aktiv werden und erneut Energie verschwenden.

\subsection{Lösungsansatz 3: Ferngesteuertes Defogging}
Die effektivste Lösung stellt ein ferngesteuertes Defogging-System dar. In diesem Szenario ist das System nicht 
dauerhaft präventiv aktiv, sondern wird gezielt durch den Fahrer gesteuert. Über eine App könnten Abfahrtszeiten 
manuell festgelegt oder wiederkehrende Routinen definiert werden. Das System würde dann wenige Minuten vor der 
geplanten Abfahrt aktiviert, um den aktuellen Zustand der Scheiben zu überprüfen und bei Bedarf geeignete 
Gegenmaßnahmen einzuleiten.

Auf diese Weise wird Energie ausschließlich dann verbraucht, wenn ein tatsächlicher Bedarf besteht. Gleichzeitig 
vereinfacht sich der erforderliche Algorithmus, da nicht mehr die langfristige Entwicklung der Scheibentemperatur 
überwacht werden muss, sondern lediglich der aktuelle Zustand relevant ist. Dieser Ansatz stellt somit einen sinnvollen
Kompromiss zwischen Funktionalität und Energieeffizienz dar und bietet die praktikabelste Lösung zur Umsetzung des 
vorgestellten Konzepts.



